\section{Every Stieltjes order and anchored parameter primitives}
\label{ag:sec:stieltjes}

We use the Laurent convention
\begin{equation}\label{ag:eq:Stieltjes-convention}
 \zeta(1+\varepsilon,x)=\frac1\varepsilon+
       \sum_{m\ge0}\frac{(-1)^m}{m!}\gamma_m(x)\varepsilon^m,
                 \qquad \gamma_0(x)=-\psi(x).
\end{equation}
Define the ordinary weighted averages
\begin{equation}\label{ag:eq:Bdefinition}
 \BB_{r,m}(a,b)=\frac{c_r}{2}\int_{-1}^1
                      (1-v^2)^r\gamma_m(a+bv)\dd v.
\end{equation}
Consequently
\[
 \HH_r(1+\varepsilon;a,b)=\frac{g_r(0)}\varepsilon+
              \sum_{m\ge0}\frac{(-1)^m}{m!}\BB_{r,m}(a,b)\varepsilon^m.
\]
This expansion is locally uniform for the spatial parameters under the
conditions of Theorem~\ref{ag:thm:hurwitz}.

\subsection{Finite complete-harmonic coefficients}
For $j\ge1$, let
\begin{equation}\label{ag:eq:homogeneous}
 E_j(\varepsilon)=\prod_{n=1}^j(1-\varepsilon/n)^{-1}
               =\sum_{\ell\ge0}h_\ell^{(j)}\varepsilon^\ell.
\end{equation}
These are complete homogeneous polynomials in $1,1/2,\ldots,1/j$.
They can be calculated from generalized harmonic numbers without
expanding a product:
\begin{equation}\label{ag:eq:hrecurrence}
 h_0^{(j)}=1,\qquad
 \ell h_\ell^{(j)}=\sum_{d=1}^\ell H_j^{(d)}h_{\ell-d}^{(j)}.
\end{equation}
For example, $h_1^{(j)}=H_j$ and
$h_2^{(j)}=((H_j)^2+H_j^{(2)})/2$.

\begin{theorem}[All-order Stieltjes endpoint formula]
\label{ag:thm:Stieltjes}
Let $D_{rj}=w_{rj}(b)/j!$, and use $\zeta^{(k)}$ for a derivative
in the spectral variable. For every $m\ge0$,
\begin{align}
 \boxed{\displaystyle
 \BB_{r,m}(a,b)=(-1)^{m+1}m!
 \sum_{j=r}^{2r}D_{rj}\sum_{k=0}^{m+1}
 \frac{h_{m+1-k}^{(j)}}{k!}
 \left\{\zeta^{(k)}(-j,a+b)-(-1)^j\zeta^{(k)}(-j,a-b)\right\}.}
 \label{ag:eq:Stieltjes-all}
\end{align}
There are only $r+1$ endpoint pairs, and the largest spectral derivative
order is $m+1$.
\end{theorem}
\begin{proof}
At $s=1+\varepsilon$,
\[
 (1-s)_{j+1}=-j!\varepsilon\prod_{n=1}^j(1-\varepsilon/n).
\]
Set
$F_j(\varepsilon)=\zeta(-j+\varepsilon,a+b)
 -(-1)^j\zeta(-j+\varepsilon,a-b)$.
Theorem~\ref{ag:thm:hurwitz} becomes
\begin{equation}\label{ag:eq:H-Laurent-master}
 \HH_r(1+\varepsilon;a,b)
 =-\frac1\varepsilon\sum_{j=r}^{2r}D_{rj}
                          E_j(\varepsilon)F_j(\varepsilon).
\end{equation}
The constant coefficient of the numerator is $-g_r(0)$.
The coefficient of $\varepsilon^{m+1}$ is a finite convolution of
\eqref{ag:eq:homogeneous} with the Taylor series of $F_j$.
Comparison with \eqref{ag:eq:Stieltjes-convention} proves
\eqref{ag:eq:Stieltjes-all}. Taking a logarithmic derivative of
\eqref{ag:eq:homogeneous} proves \eqref{ag:eq:hrecurrence}.
\end{proof}

At $m=0$ the inner expression is $F_j'(0)+H_jF_j(0)$.
At $m=1$ it is
$F_j''(0)/2+H_jF_j'(0)+((H_j)^2+H_j^{(2)})F_j(0)/2$.
Thus the harmonic counterterms are required even though every spatial
integral in \eqref{ag:eq:Bdefinition} is nonsingular. They arise from
the finite endpoint denominators, not from an arbitrary regularization.

\subsection{Spatial derivatives}
Differentiation of the compact formula gives, for $L\ge0$,
\begin{equation}\label{ag:eq:H-spatial}
 \partial_a^L\HH_r(s;a,b)=(-1)^L(s)_L\HH_r(s+L;a,b).
\end{equation}
For $L\ge1$, comparison at $s=1+\varepsilon$ yields
\begin{equation}\label{ag:eq:B-spatial}
 \partial_a^L\BB_{r,m}(a,b)
 =(-1)^{m+L}m![\varepsilon^m]
       (1+\varepsilon)_L\HH_r(1+L+\varepsilon;a,b).
\end{equation}
In particular,
\[
 \partial_a\BB_{r,m}=(-1)^{m+1}
  \left\{\partial_s^m\HH_r(2;a,b)
                   +m\partial_s^{m-1}\HH_r(2;a,b)\right\},
\]
where the second term is omitted for $m=0$. The finite product
$(1+\varepsilon)_L$ gives another explicit generalized-harmonic
coefficient expansion when desired.

\subsection{First primitives and the Gamma endpoint}
\begin{corollary}[An anchored first primitive]\label{ag:cor:primitive}
For $a,a_0>b>0$ and $m\ge0$,
\begin{equation}\label{ag:eq:first-primitive}
 \int_{a_0}^a\BB_{r,m}(\alpha,b)\dd\alpha
 =\frac{(-1)^{m+1}}{m+1}
  \left.\partial_s^{m+1}\bigl(\HH_r(s;a,b)-\HH_r(s;a_0,b)\bigr)\right|_{s=0}.
\end{equation}
For $m=0$ this is the ordinary log-Gamma identity
\begin{equation}\label{ag:eq:logGamma-primitive}
 \int_{a_0}^a\BB_{r,0}(\alpha,b)\dd\alpha
 =-\frac{c_r}{2}\int_{-1}^1(1-v^2)^r
       \log\frac{\Gamma(a+bv)}{\Gamma(a_0+bv)}\dd v.
\end{equation}
\end{corollary}
\begin{proof}
Equation \eqref{ag:eq:H-spatial} with $L=1$ gives
$\partial_a\HH_r(s)=-s\HH_r(1+s)$. Its coefficient of $s^{m+1}$ is
$(-1)^{m+1}\BB_{r,m}/m!$, proving
\eqref{ag:eq:first-primitive} after differentiation and anchoring.
Use $\zeta'(0,x)=\log\Gamma(x)-\tfrac12\log(2\pi)$
\cite{DLMFzeta} in \eqref{ag:eq:Hcompact}; the constant cancels between
the two endpoints, giving \eqref{ag:eq:logGamma-primitive}.
\end{proof}

\subsection{An explicit primitive of every depth}
The same argument can be made completely normalized at arbitrary
integration depth, without introducing unspecified constants.
For $L\ge1$ set
\[
 U_L(s;a,b)=\frac{\HH_r(s-L;a,b)}{(1-s)_L}.
\]
Define
\begin{align}
 V_L(s;a,a_0,b)={}&U_L(s;a,b)
 -\sum_{j=0}^{L-1}\frac{(a-a_0)^j}{j!}\partial_a^jU_L(s;a_0,b)\notag\\
 &-\frac{g_r(0)(a-a_0)^L}{L!(s-1)}.
 \label{ag:eq:anchored-L}
\end{align}
\begin{theorem}[All anchored Stieltjes primitives]\label{ag:thm:all-primitives}
The function $V_L$ is holomorphic near $s=1$, and
\begin{equation}\label{ag:eq:all-primitives}
 J_{L,m}(a;a_0,b)=(-1)^m m![\varepsilon^m]
                          V_L(1+\varepsilon;a,a_0,b)
\end{equation}
satisfies
$\partial_a^LJ_{L,m}=\BB_{r,m}(a,b)$ and
$\partial_a^jJ_{L,m}(a_0;a_0,b)=0$ for $0\le j<L$.
It is the unique $L$-fold primitive with those anchor conditions.
\end{theorem}
\begin{proof}
The identity $\partial_a^LU_L=\HH_r(s;a,b)$ follows from
\eqref{ag:eq:H-spatial}. The residue of $U_L$ at $s=1$ is a polynomial
of degree at most $L$ in $a$, because
$\zeta(1-L,a+bv)=-B_L(a+bv)/L$. Its $L$th derivative is $g_r(0)$.
Taylor subtraction through degree $L-1$ therefore leaves exactly the
last pole written in \eqref{ag:eq:anchored-L}; removing it proves
holomorphy. Differentiating $L$ times leaves
$\HH_r(s;a,b)-g_r(0)/(s-1)$. Coefficient comparison proves the
claimed derivative. All lower anchor derivatives vanish by construction.
The difference of any two such primitives is a polynomial of degree
below $L$ with all its anchor derivatives zero, hence vanishes.
\end{proof}
